\section{Experiments}
\label{sec:exp}

%This section evaluates CaMedPO by answering four questions: does it improve factual accuracy over alignment baselines; how each component contributes to overall performance; whether it is compatible with diverse Med-LVLM architectures; and whether it enhances clinical relevance of responses.
\subsection{Experimental Setups}
\noindent\textbf{Evaluation Datasets.}
We evaluate on the Med-VQA benchmarks VQA-RAD\cite{lau2018dataset} and SLAKE\cite{liu2021slake}, and on the radiology report benchmarks IU X-Ray\cite{demner2015preparing} and MIMIC-CXR\cite{johnson2019mimic}. We adhere to official splits and standard preprocessing to ensure strict comparability; dataset statistics are provided in Appendix \ref{supp:data}.

\begin{table*}[t]
\centering
\caption{Performance comparison on medical VQA and report generation tasks covering SLAKE, VQA-RAD, and IU-XRay datasets. For open-ended questions, we report recall (Open); for closed-ended questions, accuracy (Closed). The BLEU score denotes the average of BLEU-1/2/3/4. +SFT indicates that the model is first fine-tuned with SFT before applying the corresponding baselines. The best and second-best results are highlighted in \textbf{bold} and \underline{underlined}, respectively.}
\resizebox{\textwidth}{!}{
\begin{tabular}{lcccccccccc}
\toprule
\multirow{2}{*}{\textbf{Models}} & \multicolumn{2}{c}{\textbf{SLAKE}} & \multicolumn{2}{c}{\textbf{VQA-RAD}} & \multicolumn{3}{c}{\textbf{IU-XRay}} & \multicolumn{3}{c}{\textbf{MIMIC-CXR}} \\
\cmidrule(lr){2-3} \cmidrule(lr){4-5} \cmidrule(lr){6-8} \cmidrule(lr){9-11}
 & Open & Closed & Open & Closed & BLEU & ROUGE-L & METEOR & BLEU & ROUGE-L & METEOR \\
\midrule
LLaVA-Med v1.5 & 47.54 & 52.99 & 29.50 & 64.54 & 14.56 & 10.31 & 10.95 & 10.25 & 9.38 & 7.71 \\
\midrule
+ STLLaVA-Med & 50.48 & 54.67 & 35.00 & 68.53 & 16.10 & 10.58 & 10.51 & 11.11 & 9.29 & 7.72 \\
+ DPO & 49.84 & 53.59 & 32.00 & 66.53 & 16.08 & 12.95 & 17.13 & 11.19 & 9.45 & 7.80 \\
+ mDPO & 51.20 & 60.29 & 35.00 & 68.13 & 18.43 & 22.38 & 26.52 & 11.55 & 9.62 & 7.90 \\
+ SPPO & 52.62 & 67.58 & 38.50 & 68.93 & 21.66 & 27.30 & 32.10 & 12.32 & 10.19 & 8.83 \\
+ MMedPO & 51.67 & 58.25 & 36.50 & 67.33 & \underline{23.49} & 29.52 & \underline{34.16} & \textbf{12.85} & \underline{11.13} & \underline{10.03} \\
\rowcolor{cyan!7}
+ \textbf{CaMedPO w/ mDPO} & 51.67 & 65.07 & 38.50 & 68.53 & 22.67 & \underline{30.65} & 33.24 & 12.31 & 10.58 & 9.88 \\
\rowcolor{cyan!7}
+ \textbf{CaMedPO w/ SPPO} & \underline{53.73} & \underline{69.02} & \textbf{42.50} & \underline{69.72} & 21.46 & 27.17 & 32.23 & 12.54 & 10.56 & 9.78 \\
\rowcolor{cyan!7}
+ \textbf{CaMedPO} & \textbf{53.11} & \textbf{71.17} & \underline{40.00} & \textbf{71.31} & \textbf{23.78} & \textbf{32.68} & \textbf{35.26} & \underline{12.65} & \textbf{11.61} & \textbf{10.37} \\
\midrule
+ SFT & 52.62 & 58.01 & 31.50 & 64.14 & 22.75 & 28.86 & 33.66 & 12.39 & 10.21 & 8.75 \\
\quad+ STLLaVA-Med & 54.53 & 61.72 & 33.50 & 70.12 & 22.97 & 28.98 & 34.05 & 12.21 & 10.12 & 8.98 \\
\quad + DPO & 53.49 & 64.23 & 33.00 & 64.94 & 23.82 & 29.97 & 34.89 & 12.37 & 10.38 & 9.10 \\
\quad+ mDPO & 53.10 & 63.88 & 34.50 & 67.33 & 23.57 & 28.35 & 33.13 & 11.77 & 9.45 & 8.91 \\
\quad+ SPPO & 58.57 & 66.99 & 41.00 & 73.31 & 23.78 & 30.05 & 34.78 & 12.87 & 12.32 & 9.45 \\
\quad+ MMedPO & 56.60 & 67.82 & 37.00 & 67.73 & 24.00 & 30.13 & 35.17 & \underline{13.28} & \underline{13.22} & 10.20 \\
\rowcolor{cyan!7}
\quad+ \textbf{CaMedPO w/ mDPO} & 57.56 & 69.02 & 40.00 & 71.31 & 25.28 & 33.69 & 37.42 & 12.79 & 11.35 & 10.04 \\
\rowcolor{cyan!7}
\quad+ \textbf{CaMedPO w/ SPPO} & \underline{61.66} & \underline{70.93} & \textbf{44.00} & \underline{74.10} & \underline{27.33} & \underline{34.38} & \underline{38.11} & 12.84 & 12.60 & \underline{11.63} \\
\rowcolor{cyan!7}
\quad+ \textbf{CaMedPO} & \textbf{62.17} & \textbf{74.40} & \underline{42.50} & \textbf{74.90} & \textbf{26.23} & \textbf{34.13} & \textbf{38.25} & \textbf{14.91} & \textbf{14.58} & \textbf{12.22} \\
\bottomrule
\end{tabular}
}
\label{tab:main_results}
\end{table*}

\noindent\textbf{Implementation Details.}
We use LLaVA-Med-1.5-mistral-7B \cite{Li2023LLaVAMed} as the base model. Med-SAM \cite{ma2024segment} with clinical query guidance is used to localize the lesion area. The visual encoder is CLIP ViT-L14 at input resolution 336. %The projector adopts mlp2x\_gelu with layer selection at -2. Image inputs use padding aspect ratio. We disable image start or end tokens, as well as patch tokens. Batches are grouped by modality length, and we use distributed training with the number of processes matched to available GPUs. 
Optimization employs LoRA with rank 128 and scaling 256. Dropout is set to 0.05 and no LoRA bias is used. Training runs for three epochs with a per-device batch size of two and gradient accumulation of one. We set the learning rate to 1e-6 with cosine scheduling and a warmup ratio of 0.03. Weight decay is zero. We save at the end of each epoch, keeping a single retained checkpoint. Prompts and context formatting are kept consistent across methods. All experiments are conducted on 8 NVIDIA A100 GPUs.

\noindent\textbf{Baselines.}
We compare CaMedPO with DPO \cite{rafailov2023direct} and its variants on medical tasks, including STLLaVA-Med \cite{sun2024stllava} and MMedPO \cite{Zhu2025MMedPO}. Note that CaMedPO can be integrated with other DPO methods for improvement by modifying $\mathcal{L}_\text{DPO}$ in Eq. \ref{eq:final_loss}. We further compare our method against two preference optimization methods for natural image. Specifically, mDPO \cite{wang2024mdpo} serves as a direct preference optimization method, while SPPO \cite{wu2025selfplay} is a self-rewarding model. We integrate our causal preference optimization framework with both methods and use them as baselines to assess comparative performance. In addition, we also evaluate MMedPO and all baseline models after SFT. Further implementation details are provided in Appendix \ref{supp:baselines}.

\noindent\textbf{Evaluation Metrics.}
For VQA, we report normalized exact-match accuracy, where clinical correctness and hallucination are assessed with the Lingshu \cite{xu2025lingshu} evaluation framework MedEvalKit. For report generation, we report BLEU Score \cite{papineni2002bleu}, ROUGE-L \cite{papineni2002bleu}, and METEOR \cite{banerjee2005meteor}. 



\subsection{Main Results}






We note that the SLAKE and VQA-RAD test sets used in MMedPO differ from the official splits. For fair comparison, we therefore re-evaluate LLaVA-Med \cite{Li2023LLaVAMed}, STLLaVA-Med \cite{sun2024stllava}, DPO \cite{rafailov2023direct}, and MMedPO \cite{Zhu2025MMedPO} using the Lingshu toolkit, while official results are retained for IU-XRay and MIMIC-CXR.
CaMedPO is introduced as a causal extension of the unweighted MMedPO and is further applied to mDPO \cite{wang2024mdpo} and SPPO \cite{wu2025selfplay} to assess the generality of the proposed causal preference optimization framework.




\noindent\textbf{Comparison with Baseline Methods.}
As reported in Table \ref{tab:main_results}, CaMedPO achieves consistent improvements over all baseline preference optimization methods across both medical VQA and report-generation tasks. When compared with MMedPO \cite{Zhu2025MMedPO}, CaMedPO exhibits clear gains on multiple benchmarks. For example, on SLAKE, CaMedPO improves open-ended recall from 51.67\% to 53.11\% and closed-ended accuracy from 58.31\% to 71.16\%. On IU-XRay, CaMedPO also yields consistent gains, improving ROUGE-L from 29.52\% to 32.68\%. These representative improvements indicate that CaMedPO provides a more reliable preference signal, yielding better alignment of medical reasoning across diverse datasets.

\noindent\textbf{Comparison with Baseline Methods Enhanced by SFT.}
As shown in Table \ref{tab:main_results}, CaMedPO continues to outperform all SFT-enhanced baselines, with more improvements than in the non-SFT setting. Relative to MMedPO \cite{Zhu2025MMedPO} after supervised fine-tuning, CaMedPO achieves larger gains across both VQA and report-generation benchmarks. For example, on SLAKE-Closed, performance increases from 67.82\% to 74.40\%, and on VQA-RAD closed-ended accuracy from 67.73\% to 74.90\%. On IU-XRay, CaMedPO further improves ROUGE-L from 30.13\% to 34.13\%. These results suggest that SFT provides a stronger reference policy that amplifies the benefits of causal preference optimization. By decoupling background-related bias from the preference signal, CaMedPO is able to extract more reliable supervision from the improved SFT initialization.

\noindent\textbf{Applicability to Different Preference Optimization Algorithms.} Beyond enhancing MMedPO, the proposed causal adjustment is also applied to mDPO and SPPO to assess its generality. In both cases, the causal variants exhibit consistent improvements over their baselines. For mDPO \cite{wang2024mdpo}, CaMedPO increases SLAKE closed-ended accuracy from 60.29\% to 65.07\% and improves IU-XRay ROUGE-L from 22.38\% to 30.65\%. For SPPO, CaMedPO further raises VQA-RAD open-ended recall from 38.50\% to 41.00\% and achieves stronger report-generation results, using the second-round SPPO outputs for evaluation. These observations indicate that the proposed causal correction provides a transferable enhancement that strengthens preference optimization algorithms.



\begin{table}[t]
\centering
\caption{Ablation study results on VQA-RAD and IU-XRay. “w/” and “w/o” denote with and without, respectively. $^\text{s}$ denotes the model enhanced with SFT.}
\resizebox{\linewidth}{!}{
\begin{tabular}{lccccc}
\toprule
\multirow{2}{*}{\textbf{Methods}} & \multicolumn{2}{c}{\textbf{VQA-RAD}} & \multicolumn{3}{c}{\textbf{IU-XRay}} \\
\cmidrule(lr){2-3} \cmidrule(lr){4-6} 
 & Open & Closed & BLEU & ROUGE-L & METEOR \\
\midrule
$\text{MMedPO}^\text{s}$ & 37.00 & 67.73 & 24.00 & 30.13 & 35.17\\
$\text{CaMedPO}^\text{s}$ & 42.50 & 74.90 & 26.23 & 34.13 & 38.25 \\
\midrule
$\text{CaMedPO}^\text{s}$ w/o $\mathcal{L}_\text{DPO}$ & 29.50 & 57.77 & 14.11 & 10.65 & 10.37 \\
$\text{CaMedPO}^\text{s}$ w/ causal-weighted $\mathcal{L}_\text{DPO}$ & 37.50 & 70.52 & 24.25 & 28.40 & 34.64 \\
$\text{CaMedPO}^\text{s}$ w/ counterfactual query & 42.00 & 72.91 & 26.56 & 34.02 & 36.68 \\
\midrule
$\text{CaMedPO}^\text{s}$ w/ MedKLIP & 38.50 & 73.31 & 25.62 & 32.07 & 37.09 \\
\midrule
$\text{MMedPO}^\text{s}$ w/ background & 26.50 & 65.34 & 15.77 & 10.36 & 10.05 \\
$\text{CaMedPO}^\text{s}$ w/ background & 22.00 & 53.78 & 13.29 & 9.45 & 10.58 \\
$\text{CaMedPO}^\text{s}$ w/ composite input & 37.50 & 69.32 & 22.35 & 28.69 & 32.57 \\
\bottomrule
\end{tabular}
}
\label{tab:ablation}
\end{table}



\subsection{Ablation Study}
In Table \ref{tab:ablation}, we select SFT-enhanced MMedPO as the CaMedPO baselines to conduct abaltion studies on VQA-RAD and IU-XRay datasets. Multi-rounds training, sample ratio of causal pairs, and compatibility analysis on other LVLMs are in Appendix \ref{supp:cases}.



\noindent\textbf{Loss Function Variants.}
We evaluate several loss function variants to understand their contribution. Removing the standard DPO term leads to a clear performance drop (for example, VQA-RAD open-ended recall decreases from 42.00\% to 29.50\%), indicating that learning from factual preference pairs remains essential. Applying causal weights to the original DPO loss, where the details is in Appendix \ref{supp:cases} improves over MMedPO but remains weaker than the full model, as reflected by the IU-XRay ROUGE-L score of 28.40\% compared 
with 34.13\%. Removing the counterfactual clinical query provides only marginal improvements, thus it is not adopted in our causal preference optimization.

\noindent\textbf{Medical Visual Localization Tool.} We further examine the impact of the lesion–background decomposition module by replacing the Med-SAM \cite{ma2024segment}
localization with MedKLIP \cite{medklip}. CaMedPO$^{\text{s}}$ w/ MedKLIP leads to a moderate performance drop, where VQA-RAD open-ended recall decreases 
% from 42.50\% to 38.50\%,
by 4\%,
indicating that accurate lesion localization contributes to stronger causal contrast. Nevertheless, the model still outperforms the original MMedPO$^{\text{s}}$ baseline in all metrics, showing that even imperfect localization retains the advantage of causal preference construction.



\noindent\textbf{Reference Inputs Comparison.} We further evaluate the effect of different visual inputs used for causal preference construction.  
When replacing the original image with the background-only image, 
MMedPO$^{\text{s}}$ w/ background achieves higher performance than 
CaMedPO$^{\text{s}}$ w/ background (e.g., VQA-RAD Closed Accuracy: 65.34\% vs.\ 53.78\%), suggesting that CaMedPO effectively suppresses the background-induced bias that MMedPO continues to exploit. We additionally test using the composite image as the sole visual input. While this variant performs better than background-only inputs, its performance remains lower than the full CaMedPO model, likely because the presence of explicit background channels still introduces confounding information.  
These observations are consistent with our causal graph formulation, where background can be the biased term.

\begin{figure}[t]
    \centering
    \includegraphics[width=0.47\textwidth]{figs/cvpr_fig5.png}
    \caption{
    % Examples demonstrating the reasonable responses generated by CaMedPO.
    CaMedPO generates reasonable responses.
    }
    \label{fig:case}
    \vspace{-10pt}
\end{figure}
\section{Qualitative Analysis}
To qualitatively assess the clinical reliability of CaMedPO, we examine representative examples comparing its outputs with those produced by LLaVA-Med and MMedPO. More cases are in Appendix \ref{supp:cases}. As illustrated in Fig. \ref{fig:case}, LLaVA-Med produces an erroneous description, incorrectly suggesting bilateral pleural effusions in the report generation task. MMedPO improves factual accuracy but still omits clinically relevant details and exhibits partially incomplete reasoning. In contrast, CaMedPO provides a more comprehensive and clinically aligned report. It correctly identifies normal cardiac size, mediastinal contours, and clear lung fields, fully matching the ground-truth interpretation and achieving the highest ROUGE score.



